\documentclass[11pt,a4paper]{article}

\usepackage[T1]{fontenc}
\usepackage{lmodern}
\usepackage{microtype}
\usepackage[margin=28mm]{geometry}
\usepackage{amsmath,amssymb,amsthm,mathtools}
\usepackage{enumitem}
\usepackage{booktabs}
\usepackage[hidelinks]{hyperref}
\usepackage{fancyhdr}

\allowdisplaybreaks
\numberwithin{equation}{section}
\setlist[enumerate]{label=(\roman*),leftmargin=2.2em}

\newtheorem{theorem}{Theorem}[section]
\newtheorem{proposition}[theorem]{Proposition}
\newtheorem{lemma}[theorem]{Lemma}
\newtheorem{corollary}[theorem]{Corollary}
\theoremstyle{definition}
\newtheorem{definition}[theorem]{Definition}
\newtheorem{example}[theorem]{Example}
\newtheorem{remark}[theorem]{Remark}

\newcommand{\Q}{\mathbb Q}
\newcommand{\Z}{\mathbb Z}
\newcommand{\C}{\mathbb C}
\newcommand{\ii}{\mathrm i}
\newcommand{\N}{\operatorname N}
\newcommand{\vord}{\operatorname v}
\newcommand{\ol}[1]{\overline{#1}}
\newcommand{\A}{\mathbb A}

\pagestyle{fancy}
\fancyhf{}
\fancyhead[L]{\small Rational solutions of $y^2+z^2=x^3+1$}
\fancyhead[R]{\small Jamal Agbanwa}
\fancyfoot[C]{\thepage}
\renewcommand{\headrulewidth}{0.4pt}

\hypersetup{
  pdftitle={All Rational Solutions of y2+z2=x3+1 and the Obstruction to Polynomial Parametrization},
  pdfauthor={Jamal Agbanwa},
  pdfsubject={Rational points, Gaussian norms, and polynomial parametrization obstruction},
  pdfkeywords={Diophantine equation, Gaussian integers, norm form, polynomial parametrization, Chatelet surface}
}

\title{\bfseries All Rational Solutions of
\texorpdfstring{$y^2+z^2=x^3+1$}{y2+z2=x3+1}\\
\large and the Obstruction to an Unrestricted Polynomial Parametrization}
\author{Jamal Agbanwa\\[2pt]\small Independent Researcher}
\date{20 July 2026}

\begin{document}
\maketitle

\begin{abstract}
We study the affine surface
\[
\mathcal S:\qquad y^2+z^2=x^3+1
\]
over the rational numbers.  A natural request is a single polynomial map
$\Q^n\to\mathcal S(\Q)$, with completely unrestricted rational parameters,
whose image is all of $\mathcal S(\Q)$.  We prove that no such map exists,
regardless of the number of parameters or the degrees of the polynomials.
In fact, if nonconstant polynomials $X,Y,Z\in\Q[T_1,\ldots,T_n]$ satisfy
$Y^2+Z^2=X^3+1$, then both $X+1$ and $X^2-X+1$ are themselves sums of two
squares in the polynomial ring.  Hence such a family cannot contain even one
point with $x=2$, because $3$ is not a sum of two rational squares.  We also
show that no finite collection of polynomial families covers all rational
points.

We then give a complete prime-free descent description.  Every rational
solution is obtained from one of two norm classes, indexed by
$\varepsilon\in\{1,3\}$, and one explicit quadratic norm equation.  This
provides a rigorous all-solutions formula without placing an infinite list of
prime conditions in the statement, while also proving why the remaining
arithmetic constraint cannot be removed by a global polynomial
parametrization.
\end{abstract}

\tableofcontents

\section{The question and the precise negative answer}

Let
\begin{equation}
 \mathcal S:\qquad y^2+z^2=x^3+1.
 \label{eq:surface}
\end{equation}
For $\alpha\in\Q(\ii)$, write
\[
 \N(\alpha)=\alpha\ol\alpha.
\]
Thus $\N(a+\ii b)=a^2+b^2$.

\begin{definition}
An \emph{unrestricted polynomial parametrization} of $\mathcal S(\Q)$ means
polynomials
\[
 X,Y,Z\in\Q[T_1,\ldots,T_n]
\]
for some finite $n$ such that
\[
 Y(\mathbf T)^2+Z(\mathbf T)^2=X(\mathbf T)^3+1
\]
identically and such that every point of $\mathcal S(\Q)$ is obtained by
specializing $\mathbf T$ to a tuple in $\Q^n$.
\end{definition}

The requested object does not exist.

\begin{theorem}[Polynomial-family obstruction]
\label{thm:polynomial-obstruction}
Let $n\ge1$, and let
\[
 X,Y,Z\in\Q[T_1,\ldots,T_n]
\]
satisfy the polynomial identity
\begin{equation}
 Y^2+Z^2=X^3+1.
 \label{eq:polynomial-identity}
\end{equation}
If $X$ is nonconstant, then there exist polynomials
$A_1,A_2,B_1,B_2\in\Q[T_1,\ldots,T_n]$ such that
\begin{equation}
 X+1=A_1^2+A_2^2
 \label{eq:Xplus1-polynomial-norm}
\end{equation}
and
\begin{equation}
 X^2-X+1=B_1^2+B_2^2.
 \label{eq:quadratic-factor-polynomial-norm}
\end{equation}
Consequently, every rational specialization of a nonconstant polynomial
family satisfies that $x+1$ is a sum of two rational squares.
\end{theorem}

\begin{corollary}[No unrestricted polynomial parametrization]
\label{cor:no-global-polynomial-param}
There is no unrestricted polynomial parametrization of all rational solutions
of \eqref{eq:surface}.  More strongly, no finite collection of polynomial
maps from affine spaces over $\Q$ covers $\mathcal S(\Q)$.
\end{corollary}

The proof of Theorem~\ref{thm:polynomial-obstruction} is given in
Section~\ref{sec:polynomial-obstruction-proof}.  It is purely algebraic and
does not rely on a dimension count or on a heuristic about the apparent
complexity of the rational points.

\section{Gaussian norms over the rational numbers}
\label{sec:gaussian-norms}

We record the exact norm criterion needed later.

\begin{proposition}[Rational two-squares criterion]
\label{prop:rational-two-squares}
Let $r\in\Q_{>0}$.  The following are equivalent.
\begin{enumerate}
\item There exist $a,b\in\Q$ such that $r=a^2+b^2$.
\item For every rational prime $q\equiv3\pmod4$, the valuation
$\vord_q(r)$ is even.
\end{enumerate}
\end{proposition}

\begin{proof}
Suppose first that $r=a^2+b^2$.  After choosing a common denominator, write
\[
 r=\frac{A^2+B^2}{C^2}
\]
with $A,B,C\in\Z$ and $C\ne0$.  If $q\equiv3\pmod4$ and
$q\mid A^2+B^2$, then $q\mid A$ and $q\mid B$.  Indeed, if $q\nmid B$,
then $(AB^{-1})^2\equiv-1\pmod q$, contradicting
$(-1)^{(q-1)/2}=-1$ and Euler's criterion.  Removing the largest common
power of $q$ from $A$ and $B$ shows that $\vord_q(A^2+B^2)$ is even.
Subtracting the even integer $2\vord_q(C)$ proves that $\vord_q(r)$ is even.

Conversely, write $r=M/N$ in lowest terms, with $M,N\in\Z_{>0}$.  The parity
assumption implies that every prime $q\equiv3\pmod4$ occurs to an even
exponent separately in $M$ and in $N$.

For completeness, $\Z[\ii]$ is Euclidean for the norm
$\N(u+\ii v)=u^2+v^2$: rounding the real and imaginary parts of a Gaussian
quotient gives a remainder of norm strictly smaller than the divisor norm.
Hence $\Z[\ii]$ is a unique factorization domain.  A prime
$q\equiv3\pmod4$ remains prime in $\Z[\ii]$, while a prime
$p\equiv1\pmod4$ splits.  To see the latter directly, Wilson's theorem gives
an integer $h$ with $h^2\equiv-1\pmod p$; a nontrivial Gaussian greatest
common divisor of $p$ and $h+\ii$ then has norm $p$.  Also
$2=\N(1+\ii)$.  It follows by multiplying Gaussian factorizations that a
positive integer is a sum of two squares exactly when every prime
congruent to $3$ modulo $4$ occurs to an even exponent.  Thus
\[
 M=A^2+B^2,
 \qquad
 N=C^2+D^2
\]
for suitable integers $A,B,C,D$.  Finally,
\[
 \frac{M}{N}
 =\N\!\left(\frac{A+\ii B}{C+\ii D}\right)
 =\left(\frac{AC+BD}{N}\right)^2
  +\left(\frac{BC-AD}{N}\right)^2.
\]
This proves the converse.
\end{proof}

\begin{lemma}
\label{lem:three-not-norm}
The rational number $3$ is not a sum of two rational squares.
\end{lemma}

\begin{proof}
Suppose $3=a^2+b^2$ with $a,b\in\Q$.  Clearing denominators and cancelling a
common factor gives coprime integers $A,B,C$ satisfying
\[
 A^2+B^2=3C^2.
\]
Modulo $3$, a square is $0$ or $1$, so the left side is divisible by $3$
only if $3\mid A$ and $3\mid B$.  The displayed equation then forces
$3\mid C$, contradicting coprimality.
\end{proof}

\section{A factorization lemma in polynomial rings}
\label{sec:factorization-lemma}

Let
\[
 R=\Q[T_1,\ldots,T_n],
 \qquad
 R_{\ii}=\Q(\ii)[T_1,\ldots,T_n],
\]
and extend complex conjugation coefficientwise to $R_{\ii}$.

\begin{lemma}[Splitting coprime rational factors of a norm]
\label{lem:coprime-norm-factors}
Let $F,G\in R\setminus\{0\}$ be relatively prime in $R_{\ii}$, and suppose
that
\begin{equation}
 FG=H\ol H
 \label{eq:FG-norm}
\end{equation}
for some $H\in R_{\ii}$.  Then there exist
$c\in\Q^\times$ and $A,B\in R_{\ii}$ such that
\begin{equation}
 F=cA\ol A,
 \qquad
 G=c^{-1}B\ol B.
 \label{eq:split-norm}
\end{equation}
\end{lemma}

\begin{proof}
The ring $R_{\ii}$ is a unique factorization domain.  Factor $F$ into
irreducibles in $R_{\ii}$.  Since $F=\ol F$, conjugation permutes its
irreducible factors.

Consider first a conjugation orbit consisting of two nonassociate factors
$\pi$ and $\ol\pi$.  Their exponents in $F$ are equal.  Their contribution to
$F$ is therefore, up to a constant, the norm of a power of one member of the
pair.  Consider next an irreducible $\rho$ associated to $\ol\rho$.  Because
$F$ and $G$ are coprime, the exponent of $\rho$ in $F$ equals its exponent in
$H\ol H$.  The latter exponent is twice the exponent contributed by $H$,
so it is even.  Hence this contribution is also a norm.  Multiplying over all
orbits gives
\[
 F=cA\ol A
\]
for some $A\in R_{\ii}$ and some unit $c\in\Q(\ii)^\times$.  Both $F$ and
$A\ol A$ are fixed by conjugation, so $c=\ol c$ and therefore
$c\in\Q^\times$.

The same argument gives $G=dC\ol C$ with $d\in\Q^\times$.  From
\eqref{eq:FG-norm}, in the fraction field of $R_{\ii}$ we obtain
\[
 cd=\N\!\left(\frac{H}{AC}\right).
\]
The left side is a nonzero rational constant.  Since $\Q$ is infinite, one
may specialize the variables to rational values avoiding all zeros and poles
of $H/(AC)$.  This shows that $cd$ is the norm of an element of
$\Q(\ii)^\times$.  Choose $\lambda\in\Q(\ii)^\times$ with
$\N(\lambda)=cd$, and put $B=\lambda C$.  Then
\[
 c^{-1}B\ol B=c^{-1}(cd)C\ol C=dC\ol C=G,
\]
which proves \eqref{eq:split-norm}.
\end{proof}

\section{Proof of the polynomial obstruction}
\label{sec:polynomial-obstruction-proof}

\begin{proof}[Proof of Theorem~\ref{thm:polynomial-obstruction}]
Put
\[
 F=X+1,
 \qquad
 G=X^2-X+1,
 \qquad
 H=Y+\ii Z.
\]
Then \eqref{eq:polynomial-identity} becomes
\begin{equation}
 FG=H\ol H.
 \label{eq:norm-factorization-polynomial}
\end{equation}
Moreover,
\begin{equation}
 G-(X-2)F=3,
 \label{eq:bezout-FG}
\end{equation}
so $F$ and $G$ are relatively prime even in $R_{\ii}$.  By
Lemma~\ref{lem:coprime-norm-factors},
\begin{equation}
 F=cA\ol A,
 \qquad
 G=c^{-1}B\ol B
 \label{eq:F-G-c}
\end{equation}
for some $c\in\Q^\times$ and $A,B\in R_{\ii}$.

It remains to prove that the constant $c$ is itself a Gaussian norm.  Fix any
monomial order on $R_{\ii}$.  Let $a,b\in\Q(\ii)^\times$ be the leading
coefficients of $A$ and $B$, respectively, and let $\ell\in\Q^\times$ be the
leading coefficient of $X$.  Since $X$ is nonconstant, the leading
coefficients in \eqref{eq:F-G-c} give
\begin{equation}
 \ell=c\,a\ol a
 \label{eq:leading-one}
\end{equation}
and
\begin{equation}
 \ell^2=c^{-1}b\ol b.
 \label{eq:leading-two}
\end{equation}
Combining these identities yields
\[
 b\ol b=c^3(a\ol a)^2.
\]
Therefore
\begin{equation}
 c=
 \N\!\left(\frac{b}{c a^2}\right).
 \label{eq:c-is-norm}
\end{equation}
Set $\delta=b/(ca^2)\in\Q(\ii)^\times$.  Then $\N(\delta)=c$, and
\eqref{eq:F-G-c} becomes
\[
 F=(\delta A)\ol{(\delta A)},
 \qquad
 G=(B/\delta)\ol{(B/\delta)}.
\]
Writing
\[
 \delta A=A_1+\ii A_2,
 \qquad
 B/\delta=B_1+\ii B_2
\]
with $A_1,A_2,B_1,B_2\in R$ proves
\eqref{eq:Xplus1-polynomial-norm} and
\eqref{eq:quadratic-factor-polynomial-norm}.
\end{proof}

\begin{proof}[Proof of Corollary~\ref{cor:no-global-polynomial-param}]
The surface contains the rational point $(2,3,0)$.  If a nonconstant
polynomial family contained this point, Theorem~\ref{thm:polynomial-obstruction}
would imply that
\[
 3=2+1
\]
is a sum of two rational squares, contradicting
Lemma~\ref{lem:three-not-norm}.  Hence a single surjective polynomial family
cannot exist.

For the stronger finite-cover assertion, first note that if $X$ is constant
in a polynomial family, then
\[
 (Y+\ii Z)(Y-\ii Z)=X^3+1
\]
is constant.  In the polynomial ring $R_{\ii}$, both factors on the left are
therefore units, so $Y$ and $Z$ are constant as well.  Thus a polynomial
family with constant $X$ has a one-point image.

On the other hand, the fiber $x=2$ contains infinitely many rational points:
for every $t\in\Q$,
\begin{equation}
 \left(
 2,
 3\frac{1-t^2}{1+t^2},
 \frac{6t}{1+t^2}
 \right)
 \in\mathcal S(\Q).
 \label{eq:x2-fiber}
\end{equation}
No nonconstant polynomial family meets this fiber, while each constant
polynomial family contributes at most one point.  A finite collection can
therefore cover only finitely many of the points in \eqref{eq:x2-fiber}.
\end{proof}

\begin{remark}
Theorem~\ref{thm:polynomial-obstruction} is stronger than the elementary
observation that one particular proposed formula fails.  It rules out every
possible polynomial formula, with any finite number of unrestricted rational
parameters and of any degree.
\end{remark}

\section{A complete prime-free two-class descent}
\label{sec:two-class-descent}

The impossibility theorem does not prevent an exact all-solutions
description.  It shows instead that such a description must retain genuine
arithmetic data.  For this equation the data can be compressed into the two
classes $1$ and $3$.

\begin{lemma}[The two possible norm classes]
\label{lem:two-norm-classes}
Let $x\in\Q$ with $x>-1$, and suppose that $x^3+1$ is a sum of two rational
squares.  Then there is a unique
\[
 \varepsilon\in\{1,3\}
\]
such that both
\begin{equation}
 \frac{x+1}{\varepsilon}
 \qquad\text{and}\qquad
 \frac{x^2-x+1}{\varepsilon}
 \label{eq:two-factor-norms}
\end{equation}
are sums of two rational squares.
\end{lemma}

\begin{proof}
Write $x=a/b$ in lowest terms with $b>0$, and set
\[
 U=x+1=\frac{a+b}{b},
 \qquad
 V=x^2-x+1=\frac{a^2-ab+b^2}{b^2}.
\]
Then $UV=x^3+1$ is a rational Gaussian norm.

Let $q\equiv3\pmod4$ be a prime with $q\ne3$.  If $q\mid b$, then
$q\nmid a$, and hence
\[
 \vord_q(U)=-\vord_q(b),
 \qquad
 \vord_q(V)=-2\vord_q(b).
\]
Since $\vord_q(UV)$ is even, $\vord_q(b)$ is even, so both displayed
valuations are even.

Suppose now that $q\nmid b$.  If $q$ divided both $a+b$ and
$a^2-ab+b^2$, then $a\equiv-b\pmod q$ and consequently
\[
 a^2-ab+b^2\equiv3b^2\pmod q,
\]
forcing $q=3$, a contradiction.  Thus at most one of $U$ and $V$ has
positive $q$-adic valuation.  Since their valuation sum is even, each
individual valuation is even.

It follows that, for every prime $q\equiv3\pmod4$ other than $3$, the
$q$-adic valuations of both $U$ and $V$ are even.  At $q=3$, the parity of
$\vord_3(U)+\vord_3(V)$ is even, so $\vord_3(U)$ and $\vord_3(V)$ have the
same parity.  Choose $\varepsilon=1$ when this common parity is even, and
$\varepsilon=3$ when it is odd.  Proposition~\ref{prop:rational-two-squares}
then shows that both numbers in \eqref{eq:two-factor-norms} are sums of two
rational squares.

Uniqueness follows from Lemma~\ref{lem:three-not-norm}: if both choices of
$\varepsilon$ worked, their quotient would imply that $3$ is a rational
Gaussian norm.
\end{proof}

\begin{theorem}[Complete prime-free rational classification]
\label{thm:prime-free-classification}
A triple $(x,y,z)\in\Q^3$ satisfies \eqref{eq:surface} if and only if either
\[
 (x,y,z)=(-1,0,0),
\]
or there exist
\[
 \varepsilon\in\{1,3\},
 \qquad
 a,b,c,d,m,n\in\Q,
 \qquad
 m^2+n^2\ne0,
\]
such that, with
\begin{equation}
 S=a^2+b^2,
 \label{eq:def-S}
\end{equation}
one has
\begin{equation}
 c^2+d^2
 =\varepsilon S^2-3S+\frac{3}{\varepsilon},
 \label{eq:single-residual-norm-equation}
\end{equation}
and
\begin{equation}
 \boxed{
 x=\varepsilon S-1,
 \qquad
 y+\ii z=
 \varepsilon(a+\ii b)(c+\ii d)
 \frac{m+\ii n}{m-\ii n}.
 }
 \label{eq:prime-free-gaussian-formula}
\end{equation}
The boundary solution $(-1,0,0)$ may also be included in the formula by
taking $\varepsilon=3$, $a=b=0$, and $c^2+d^2=1$.
\end{theorem}

\begin{proof}
Assume first that $(x,y,z)$ is a rational solution with $x>-1$.  Put
$W=y+\ii z$.  By Lemma~\ref{lem:two-norm-classes}, choose the unique
$\varepsilon\in\{1,3\}$ and Gaussian rationals
\[
 \alpha=a+\ii b,
 \qquad
 \beta=c+\ii d
\]
such that
\begin{equation}
 \N(\alpha)=\frac{x+1}{\varepsilon},
 \qquad
 \N(\beta)=\frac{x^2-x+1}{\varepsilon}.
 \label{eq:alpha-beta-norms}
\end{equation}
The first equation gives $x=\varepsilon S-1$.  Substitution into the second
gives
\[
 c^2+d^2
 =\frac{(\varepsilon S-1)^2-(\varepsilon S-1)+1}{\varepsilon}
 =\varepsilon S^2-3S+\frac3\varepsilon,
\]
which is \eqref{eq:single-residual-norm-equation}.

Furthermore,
\[
 \N(\varepsilon\alpha\beta)
 =\varepsilon^2\N(\alpha)\N(\beta)
 =(x+1)(x^2-x+1)
 =x^3+1
 =\N(W).
\]
Hence
\[
 \eta=\frac{W}{\varepsilon\alpha\beta}
\]
has norm $1$.  Every norm-one Gaussian rational has the form
\begin{equation}
 \eta=\frac{m+\ii n}{m-\ii n}
 \label{eq:hilbert90-param}
\end{equation}
with $m,n\in\Q$ not both zero.  Indeed, if
$\eta=C+\ii D$ and $C\ne-1$, take $m=1+C$ and $n=D$; if $\eta=-1$, take
$m=0$ and $n=1$.  This proves \eqref{eq:prime-free-gaussian-formula}.

Conversely, suppose the parameters satisfy
\eqref{eq:single-residual-norm-equation} and define $x,y,z$ by
\eqref{eq:prime-free-gaussian-formula}.  Since the last factor has norm $1$,
\begin{align*}
 y^2+z^2
 &=\varepsilon^2S(c^2+d^2)\\
 &=\varepsilon^2S\left(\varepsilon S^2-3S+\frac3\varepsilon\right)\\
 &=(\varepsilon S-1)^3+1\\
 &=x^3+1.
\end{align*}
Thus every permitted parameter choice gives a rational solution, and the
preceding argument proves that every rational solution occurs.
\end{proof}

\begin{corollary}[Expanded real-coordinate formula]
\label{cor:expanded-real}
Under the hypotheses of Theorem~\ref{thm:prime-free-classification}, put
\[
 P=ac-bd,
 \qquad
 Q=ad+bc.
\]
Then
\begin{equation}
 \boxed{
 y=\varepsilon\,
 \frac{(m^2-n^2)P-2mnQ}{m^2+n^2},
 }
 \label{eq:y-expanded}
\end{equation}
\begin{equation}
 \boxed{
 z=\varepsilon\,
 \frac{2mnP+(m^2-n^2)Q}{m^2+n^2}.
 }
 \label{eq:z-expanded}
\end{equation}
Together with
\[
 x=\varepsilon(a^2+b^2)-1
\]
and \eqref{eq:single-residual-norm-equation}, these formulas give every
rational solution.
\end{corollary}

\begin{remark}[What remains constrained]
Theorem~\ref{thm:prime-free-classification} removes the infinite prime test
from the statement and replaces it by the finite choice
$\varepsilon\in\{1,3\}$ and the single explicit equation
\eqref{eq:single-residual-norm-equation}.  Theorem~\ref{thm:polynomial-obstruction}
shows that one cannot eliminate all arithmetic constraints and still obtain
all rational points from one unrestricted polynomial family.
\end{remark}

\section{Examples and explicit unrestricted rational families}

\begin{example}[The principal norm class]
Take $\varepsilon=1$ and $t\in\Q$.  The identity
\[
 1^2+(t^3)^2=(t^2)^3+1
\]
gives the polynomial family
\begin{equation}
 (x,y,z)=(t^2,1,t^3).
 \label{eq:square-x-family}
\end{equation}
Here $x+1=t^2+1$ is visibly a rational Gaussian norm, in accordance with
Theorem~\ref{thm:polynomial-obstruction}.
\end{example}

\begin{example}[The nonprincipal norm class]
The point
\[
 \left(\frac12,\frac34,\frac34\right)
\]
is a rational solution because
\[
 \frac{9}{16}+\frac{9}{16}=\frac98
 =\left(\frac12\right)^3+1.
\]
It belongs to the $\varepsilon=3$ branch.  Indeed,
\[
 \frac{x+1}{3}=\frac12
 =\N\!\left(\frac{1+\ii}{2}\right),
 \qquad
 \frac{x^2-x+1}{3}=\frac14
 =\N\!\left(\frac12\right),
\]
and
\[
 \frac34+\frac34\ii
 =3\left(\frac{1+\ii}{2}\right)\left(\frac12\right).
\]
No nonconstant polynomial family over $\Q$ can contain this point, because
$x+1=3/2$ is not a rational Gaussian norm.
\end{example}

\begin{example}[Why denominators matter]
Although the $\varepsilon=3$ branch cannot occur in a nonconstant polynomial
family, it does occur in unrestricted rational-function families.  Let
$t\in\Q$.  Since $t^2-2\ne0$ for every rational $t$, the formulas
\begin{equation}
 \begin{aligned}
 x(t)&=2+\frac{12(t^2-1)}{(t^2-2)^2},\\
 y(t)&=3-\frac{24(t^2-1)}{(t^2-2)^3},\\
 z(t)&=\frac{12t^3(t^2-1)}{(t^2-2)^3}
 \end{aligned}
 \label{eq:class-three-rational-family}
\end{equation}
are defined for every $t\in\Q$ and satisfy
\[
 y(t)^2+z(t)^2=x(t)^3+1.
\]
Moreover,
\begin{equation}
 x(t)+1=3\left(\frac{t^2}{t^2-2}\right)^2,
 \label{eq:class-three-norm-class}
\end{equation}
so the generic member lies in the $\varepsilon=3$ class.

For completeness, the family is obtained by taking the line through
\[
 P=(2,3,0)
 \qquad\text{and}\qquad
 Q_t=(t^2,1,t^3).
\]
Writing the line as $P+\lambda(Q_t-P)$ and substituting into
$y^2+z^2-x^3-1$ gives
\[
 -\lambda(\lambda-1)
 \left((t^2-2)^3\lambda-12(t^2-1)\right).
\]
The third intersection parameter is therefore
\[
 \lambda=\frac{12(t^2-1)}{(t^2-2)^3},
\]
which yields \eqref{eq:class-three-rational-family}.
\end{example}

An arbitrary rational rotation of either family gives a second unrestricted
parameter.  If $(X,Y,Z)$ is any rational solution and $u\in\Q$, set
\begin{equation}
 \begin{aligned}
 Y_u&=\frac{(1-u^2)Y-2uZ}{1+u^2},\\
 Z_u&=\frac{2uY+(1-u^2)Z}{1+u^2}.
 \end{aligned}
 \label{eq:rational-rotation}
\end{equation}
Then $Y_u^2+Z_u^2=Y^2+Z^2$.  Since $1+u^2\ne0$ for rational $u$, the
parameter is unrestricted.  These rational-function families are useful and
Zariski-dense, but they are not a complete parametrization of all rational
points.

\section{Integral points}

For completeness, the exact integral criterion is immediate from
Proposition~\ref{prop:rational-two-squares}.

\begin{corollary}
Let $x,y,z\in\Z$.  Then $y^2+z^2=x^3+1$ if and only if either
$(x,y,z)=(-1,0,0)$, or $x\ge0$, every prime $q\equiv3\pmod4$ occurs to an
even exponent in $x^3+1$, and $(y,z)$ is one of the integral two-square
representations of $x^3+1$.
\end{corollary}

\begin{proof}
The inequality $y^2+z^2\ge0$ gives $x\ge-1$.  The case $x=-1$ is immediate.
For $x\ge0$, Proposition~\ref{prop:rational-two-squares}, applied to the
integer $x^3+1$, gives the criterion.  Every representation of that integer
as $y^2+z^2$ is plainly a solution, and there are no others.
\end{proof}

\section{Conclusion}

The demand for one explicit polynomial parametrization with unrestricted
rational parameters is incompatible with the arithmetic of
\eqref{eq:surface}.  It is not merely that a convenient formula has not yet
been found.  Theorem~\ref{thm:polynomial-obstruction} proves that every
nonconstant polynomial family is trapped in the norm class in which $x+1$ is
a sum of two rational squares, whereas the surface has infinitely many
rational points in the other class, already on the fiber $x=2$.

The complete replacement is Theorem~\ref{thm:prime-free-classification}:
all rational solutions are described by the finite branch
$\varepsilon\in\{1,3\}$, one explicit quadratic norm equation, and an
unrestricted norm-one rotation.  This is both exact and structurally sharp.

\end{document}
